\documentclass[10pt]{article}
% Hojas resumen: un A4 por delante y por detrás, dos columnas por cara.
% Uso: \documentclass[10pt]{article} % Hojas resumen: un A4 por delante y por detrás, dos columnas por cara.
% Uso: \documentclass[10pt]{article} % Hojas resumen: un A4 por delante y por detrás, dos columnas por cara.
% Uso: \documentclass[10pt]{article} \input{estilo-resumen.tex} \bloque{analisis|algebra|azar} \cabecera{n}{Título}{Lecciones}
\usepackage[utf8]{inputenc}
\usepackage[T1]{fontenc}
\usepackage[spanish,es-nodecimaldot]{babel}
\usepackage{lmodern}
\usepackage{amsmath,amssymb}
\usepackage{graphicx}
\usepackage[a4paper,left=10mm,right=10mm,top=15mm,bottom=13mm,headheight=11pt,headsep=3mm,footskip=6mm]{geometry}
\usepackage{multicol}
\setlength{\columnsep}{6mm}
\setlength{\columnseprule}{0.3pt}
\usepackage{xcolor}
\usepackage{enumitem}
\setlist{nosep,leftmargin=1.1em,itemsep=1pt}
\usepackage{array}
\usepackage{needspace}
\usepackage{pgfplots}
\usepgfplotslibrary{fillbetween}
\pgfplotsset{compat=1.18,
  mini/.style={width=4.3cm,height=3.3cm,scale only axis=false,axis lines=middle,xtick=\empty,ytick=\empty,samples=120,
    xlabel={\scriptsize$x$},ylabel={\scriptsize$y$},title style={font=\scriptsize,yshift=-4pt},clip=true,
    every axis plot/.append style={thick,color=acento}}}
\usepackage{ragged2e}
\usepackage[most]{tcolorbox}
\usepackage{fancyhdr,lastpage}
\usepackage[hidelinks,pdfencoding=auto]{hyperref}
\hypersetup{pdflang=es-ES,pdfsubject={Matemáticas Aplicadas a las Ciencias Sociales II}}

\definecolor{analisis}{HTML}{1D4ED8}
\definecolor{algebra}{HTML}{86198F}
\definecolor{azar}{HTML}{9A3412}
\definecolor{aviso}{HTML}{B91C1C}
\colorlet{acento}{analisis}
\newcommand{\bloque}[1]{\colorlet{acento}{#1}}

\setlength{\parindent}{0pt}
\setlength{\parskip}{2pt}
\renewcommand{\arraystretch}{1.2}
\AtBeginDocument{\RaggedRight\fontsize{9.8}{11.9}\selectfont\setlength{\abovedisplayskip}{3pt}\setlength{\belowdisplayskip}{3pt}\setlength{\abovedisplayshortskip}{1pt}\setlength{\belowdisplayshortskip}{1pt}}

\newcommand{\temanum}{}\newcommand{\tematitulo}{}
\pagestyle{fancy}
\fancyhf{}
\fancyhead[L]{\footnotesize\color{acento}\textbf{Hoja resumen}\quad Tema \temanum. \tematitulo}
\fancyhead[R]{\footnotesize Matemáticas Aplicadas a las Ciencias Sociales II --- 2.\textsuperscript{o} Bachillerato}
\fancyfoot[L]{\scriptsize IES Al-Ándalus --- Juan María Gámez Ortiz, con la ayuda de Claude y de Codex --- CC BY-NC 4.0}
\fancyfoot[R]{\scriptsize Cara \thepage\ de \pageref{LastPage}}
\renewcommand{\headrulewidth}{0.4pt}

% Cabecera de la primera cara: número, título y lecciones del tema.
\newcommand{\cabecera}[3]{%
  \renewcommand{\temanum}{#1}\renewcommand{\tematitulo}{#2}%
  \hypersetup{pdftitle={Hoja resumen. Tema #1. #2}}%
  \noindent\begin{tcolorbox}[enhanced,colback=acento!8,colframe=acento,boxrule=0.6pt,arc=2pt,left=4pt,right=4pt,top=2pt,bottom=2pt]
  \begin{minipage}[c]{1.1cm}\includegraphics[width=1cm]{logo-al-andalus.png}\end{minipage}%
  \begin{minipage}[c]{\dimexpr\linewidth-1.2cm}{\large\bfseries\color{acento}Tema #1. #2}\\[1pt]{\scriptsize #3}\end{minipage}
  \end{tcolorbox}\vspace{-2pt}}

% Apartado con filete del color del bloque.
\newcommand{\seccion}[1]{\par\needspace{5\baselineskip}\vspace{3pt}{\color{acento}\bfseries\normalsize #1}\par\vspace{-1pt}{\color{acento}\hrule height 0.5pt}\vspace{2pt}}
\newcommand{\sub}[1]{\par\needspace{3\baselineskip}\vspace{1pt}\textbf{\color{acento}#1}\enspace}

% Fórmula clave, aviso de errores y ejemplo.
\tcbset{before upper=\RaggedRight}
\newtcolorbox{formula}{enhanced,colback=acento!6,colframe=acento!70,boxrule=0.5pt,arc=1.5pt,left=3pt,right=3pt,top=0pt,bottom=0pt,before skip=2pt,after skip=2pt}
\newtcolorbox{ojo}{enhanced,colback=aviso!5,colframe=aviso!60,boxrule=0.5pt,arc=1.5pt,left=3pt,right=3pt,top=1pt,bottom=1pt,before skip=3pt,after skip=2pt,
  fonttitle=\bfseries\footnotesize,coltitle=aviso,colbacktitle=aviso!5,title={Ojo},attach title to upper={\enspace}}
\newtcolorbox{ejemplo}{enhanced,colback=white,colframe=acento!45,boxrule=0.4pt,arc=1.5pt,left=3pt,right=3pt,top=1pt,bottom=1pt,before skip=3pt,after skip=2pt,
  fonttitle=\bfseries\footnotesize,coltitle=acento,colbacktitle=white,title={Ejemplo},attach title to upper={\enspace}}
\newcommand{\dx}{\,dx}
 \bloque{analisis|algebra|azar} \cabecera{n}{Título}{Lecciones}
\usepackage[utf8]{inputenc}
\usepackage[T1]{fontenc}
\usepackage[spanish,es-nodecimaldot]{babel}
\usepackage{lmodern}
\usepackage{amsmath,amssymb}
\usepackage{graphicx}
\usepackage[a4paper,left=10mm,right=10mm,top=15mm,bottom=13mm,headheight=11pt,headsep=3mm,footskip=6mm]{geometry}
\usepackage{multicol}
\setlength{\columnsep}{6mm}
\setlength{\columnseprule}{0.3pt}
\usepackage{xcolor}
\usepackage{enumitem}
\setlist{nosep,leftmargin=1.1em,itemsep=1pt}
\usepackage{array}
\usepackage{needspace}
\usepackage{pgfplots}
\usepgfplotslibrary{fillbetween}
\pgfplotsset{compat=1.18,
  mini/.style={width=4.3cm,height=3.3cm,scale only axis=false,axis lines=middle,xtick=\empty,ytick=\empty,samples=120,
    xlabel={\scriptsize$x$},ylabel={\scriptsize$y$},title style={font=\scriptsize,yshift=-4pt},clip=true,
    every axis plot/.append style={thick,color=acento}}}
\usepackage{ragged2e}
\usepackage[most]{tcolorbox}
\usepackage{fancyhdr,lastpage}
\usepackage[hidelinks,pdfencoding=auto]{hyperref}
\hypersetup{pdflang=es-ES,pdfsubject={Matemáticas Aplicadas a las Ciencias Sociales II}}

\definecolor{analisis}{HTML}{1D4ED8}
\definecolor{algebra}{HTML}{86198F}
\definecolor{azar}{HTML}{9A3412}
\definecolor{aviso}{HTML}{B91C1C}
\colorlet{acento}{analisis}
\newcommand{\bloque}[1]{\colorlet{acento}{#1}}

\setlength{\parindent}{0pt}
\setlength{\parskip}{2pt}
\renewcommand{\arraystretch}{1.2}
\AtBeginDocument{\RaggedRight\fontsize{9.8}{11.9}\selectfont\setlength{\abovedisplayskip}{3pt}\setlength{\belowdisplayskip}{3pt}\setlength{\abovedisplayshortskip}{1pt}\setlength{\belowdisplayshortskip}{1pt}}

\newcommand{\temanum}{}\newcommand{\tematitulo}{}
\pagestyle{fancy}
\fancyhf{}
\fancyhead[L]{\footnotesize\color{acento}\textbf{Hoja resumen}\quad Tema \temanum. \tematitulo}
\fancyhead[R]{\footnotesize Matemáticas Aplicadas a las Ciencias Sociales II --- 2.\textsuperscript{o} Bachillerato}
\fancyfoot[L]{\scriptsize IES Al-Ándalus --- Juan María Gámez Ortiz, con la ayuda de Claude y de Codex --- CC BY-NC 4.0}
\fancyfoot[R]{\scriptsize Cara \thepage\ de \pageref{LastPage}}
\renewcommand{\headrulewidth}{0.4pt}

% Cabecera de la primera cara: número, título y lecciones del tema.
\newcommand{\cabecera}[3]{%
  \renewcommand{\temanum}{#1}\renewcommand{\tematitulo}{#2}%
  \hypersetup{pdftitle={Hoja resumen. Tema #1. #2}}%
  \noindent\begin{tcolorbox}[enhanced,colback=acento!8,colframe=acento,boxrule=0.6pt,arc=2pt,left=4pt,right=4pt,top=2pt,bottom=2pt]
  \begin{minipage}[c]{1.1cm}\includegraphics[width=1cm]{logo-al-andalus.png}\end{minipage}%
  \begin{minipage}[c]{\dimexpr\linewidth-1.2cm}{\large\bfseries\color{acento}Tema #1. #2}\\[1pt]{\scriptsize #3}\end{minipage}
  \end{tcolorbox}\vspace{-2pt}}

% Apartado con filete del color del bloque.
\newcommand{\seccion}[1]{\par\needspace{5\baselineskip}\vspace{3pt}{\color{acento}\bfseries\normalsize #1}\par\vspace{-1pt}{\color{acento}\hrule height 0.5pt}\vspace{2pt}}
\newcommand{\sub}[1]{\par\needspace{3\baselineskip}\vspace{1pt}\textbf{\color{acento}#1}\enspace}

% Fórmula clave, aviso de errores y ejemplo.
\tcbset{before upper=\RaggedRight}
\newtcolorbox{formula}{enhanced,colback=acento!6,colframe=acento!70,boxrule=0.5pt,arc=1.5pt,left=3pt,right=3pt,top=0pt,bottom=0pt,before skip=2pt,after skip=2pt}
\newtcolorbox{ojo}{enhanced,colback=aviso!5,colframe=aviso!60,boxrule=0.5pt,arc=1.5pt,left=3pt,right=3pt,top=1pt,bottom=1pt,before skip=3pt,after skip=2pt,
  fonttitle=\bfseries\footnotesize,coltitle=aviso,colbacktitle=aviso!5,title={Ojo},attach title to upper={\enspace}}
\newtcolorbox{ejemplo}{enhanced,colback=white,colframe=acento!45,boxrule=0.4pt,arc=1.5pt,left=3pt,right=3pt,top=1pt,bottom=1pt,before skip=3pt,after skip=2pt,
  fonttitle=\bfseries\footnotesize,coltitle=acento,colbacktitle=white,title={Ejemplo},attach title to upper={\enspace}}
\newcommand{\dx}{\,dx}
 \bloque{analisis|algebra|azar} \cabecera{n}{Título}{Lecciones}
\usepackage[utf8]{inputenc}
\usepackage[T1]{fontenc}
\usepackage[spanish,es-nodecimaldot]{babel}
\usepackage{lmodern}
\usepackage{amsmath,amssymb}
\usepackage{graphicx}
\usepackage[a4paper,left=10mm,right=10mm,top=15mm,bottom=13mm,headheight=11pt,headsep=3mm,footskip=6mm]{geometry}
\usepackage{multicol}
\setlength{\columnsep}{6mm}
\setlength{\columnseprule}{0.3pt}
\usepackage{xcolor}
\usepackage{enumitem}
\setlist{nosep,leftmargin=1.1em,itemsep=1pt}
\usepackage{array}
\usepackage{needspace}
\usepackage{pgfplots}
\usepgfplotslibrary{fillbetween}
\pgfplotsset{compat=1.18,
  mini/.style={width=4.3cm,height=3.3cm,scale only axis=false,axis lines=middle,xtick=\empty,ytick=\empty,samples=120,
    xlabel={\scriptsize$x$},ylabel={\scriptsize$y$},title style={font=\scriptsize,yshift=-4pt},clip=true,
    every axis plot/.append style={thick,color=acento}}}
\usepackage{ragged2e}
\usepackage[most]{tcolorbox}
\usepackage{fancyhdr,lastpage}
\usepackage[hidelinks,pdfencoding=auto]{hyperref}
\hypersetup{pdflang=es-ES,pdfsubject={Matemáticas Aplicadas a las Ciencias Sociales II}}

\definecolor{analisis}{HTML}{1D4ED8}
\definecolor{algebra}{HTML}{86198F}
\definecolor{azar}{HTML}{9A3412}
\definecolor{aviso}{HTML}{B91C1C}
\colorlet{acento}{analisis}
\newcommand{\bloque}[1]{\colorlet{acento}{#1}}

\setlength{\parindent}{0pt}
\setlength{\parskip}{2pt}
\renewcommand{\arraystretch}{1.2}
\AtBeginDocument{\RaggedRight\fontsize{9.8}{11.9}\selectfont\setlength{\abovedisplayskip}{3pt}\setlength{\belowdisplayskip}{3pt}\setlength{\abovedisplayshortskip}{1pt}\setlength{\belowdisplayshortskip}{1pt}}

\newcommand{\temanum}{}\newcommand{\tematitulo}{}
\pagestyle{fancy}
\fancyhf{}
\fancyhead[L]{\footnotesize\color{acento}\textbf{Hoja resumen}\quad Tema \temanum. \tematitulo}
\fancyhead[R]{\footnotesize Matemáticas Aplicadas a las Ciencias Sociales II --- 2.\textsuperscript{o} Bachillerato}
\fancyfoot[L]{\scriptsize IES Al-Ándalus --- Juan María Gámez Ortiz, con la ayuda de Claude y de Codex --- CC BY-NC 4.0}
\fancyfoot[R]{\scriptsize Cara \thepage\ de \pageref{LastPage}}
\renewcommand{\headrulewidth}{0.4pt}

% Cabecera de la primera cara: número, título y lecciones del tema.
\newcommand{\cabecera}[3]{%
  \renewcommand{\temanum}{#1}\renewcommand{\tematitulo}{#2}%
  \hypersetup{pdftitle={Hoja resumen. Tema #1. #2}}%
  \noindent\begin{tcolorbox}[enhanced,colback=acento!8,colframe=acento,boxrule=0.6pt,arc=2pt,left=4pt,right=4pt,top=2pt,bottom=2pt]
  \begin{minipage}[c]{1.1cm}\includegraphics[width=1cm]{logo-al-andalus.png}\end{minipage}%
  \begin{minipage}[c]{\dimexpr\linewidth-1.2cm}{\large\bfseries\color{acento}Tema #1. #2}\\[1pt]{\scriptsize #3}\end{minipage}
  \end{tcolorbox}\vspace{-2pt}}

% Apartado con filete del color del bloque.
\newcommand{\seccion}[1]{\par\needspace{5\baselineskip}\vspace{3pt}{\color{acento}\bfseries\normalsize #1}\par\vspace{-1pt}{\color{acento}\hrule height 0.5pt}\vspace{2pt}}
\newcommand{\sub}[1]{\par\needspace{3\baselineskip}\vspace{1pt}\textbf{\color{acento}#1}\enspace}

% Fórmula clave, aviso de errores y ejemplo.
\tcbset{before upper=\RaggedRight}
\newtcolorbox{formula}{enhanced,colback=acento!6,colframe=acento!70,boxrule=0.5pt,arc=1.5pt,left=3pt,right=3pt,top=0pt,bottom=0pt,before skip=2pt,after skip=2pt}
\newtcolorbox{ojo}{enhanced,colback=aviso!5,colframe=aviso!60,boxrule=0.5pt,arc=1.5pt,left=3pt,right=3pt,top=1pt,bottom=1pt,before skip=3pt,after skip=2pt,
  fonttitle=\bfseries\footnotesize,coltitle=aviso,colbacktitle=aviso!5,title={Ojo},attach title to upper={\enspace}}
\newtcolorbox{ejemplo}{enhanced,colback=white,colframe=acento!45,boxrule=0.4pt,arc=1.5pt,left=3pt,right=3pt,top=1pt,bottom=1pt,before skip=3pt,after skip=2pt,
  fonttitle=\bfseries\footnotesize,coltitle=acento,colbacktitle=white,title={Ejemplo},attach title to upper={\enspace}}
\newcommand{\dx}{\,dx}

\bloque{azar}
\begin{document}
\cabecera{7}{Probabilidad}{Lecciones 7.1--7.10: sucesos y regla de Laplace, álgebra de sucesos, tablas de contingencia, frecuencias, axiomas de Kolmogórov, propiedades, experimentos compuestos, probabilidad condicionada, probabilidad total y teorema de Bayes.}
\begin{multicols}{2}

\seccion{1. Experimentos y sucesos}
\begin{itemize}
  \item \textbf{Experimento aleatorio}: se conocen los resultados posibles, pero no cuál saldrá.
  \item \textbf{Espacio muestral} \(\Omega\): todos los resultados. \textbf{Suceso}: subconjunto de \(\Omega\).
  \item \textbf{Elemental} (un resultado), \textbf{seguro} (\(\Omega\)), \textbf{imposible} (\(\varnothing\)).
\end{itemize}
\begin{formula}
\textbf{Regla de Laplace} (resultados \textbf{equiprobables}):
\[P(A)=\frac{|A|}{|\Omega|}=\frac{\text{casos favorables}}{\text{casos posibles}}\]
\end{formula}
\begin{ejemplo}
Un dado: \(P(\text{par})=\frac36=\frac12\). Dos dados: \(P(\text{suma }7)=\frac6{36}=\frac16\).
\end{ejemplo}
\begin{ojo}
Laplace solo vale si los resultados son equiprobables: en una ruleta con sectores desiguales, no.
\end{ojo}

\seccion{2. Álgebra de sucesos}
\begin{center}
\begin{tikzpicture}[scale=0.62]
\fill[acento!18] (0,0) circle (1.2);
\fill[acento!18] (1.5,0) circle (1.2);
\begin{scope}\clip (0,0) circle (1.2); \fill[acento!50] (1.5,0) circle (1.2);\end{scope}
\draw[thick,acento] (0,0) circle (1.2) (1.5,0) circle (1.2);
\draw (-1.8,-1.6) rectangle (3.3,1.6);
\node at (-0.5,0) {\small \(A\)}; \node at (2,0) {\small \(B\)}; \node at (0.75,0) {\scriptsize \(A\!\cap\! B\)};
\node at (2.9,1.25) {\small \(\Omega\)};
\end{tikzpicture}
\end{center}
\begin{itemize}
  \item \textbf{Unión} \(A\cup B\): ocurre \(A\) o \(B\) (o ambos). \textbf{Intersección} \(A\cap B\): ocurren los dos.
  \item \textbf{Contrario} \(A^c\): no ocurre \(A\). \textbf{Diferencia} \(A-B=A\cap B^c\).
  \item \textbf{Incompatibles}: \(A\cap B=\varnothing\) (no pueden ocurrir a la vez).
\end{itemize}
\begin{formula}
\textbf{Leyes de De Morgan:}
\[(A\cup B)^c=A^c\cap B^c,\qquad (A\cap B)^c=A^c\cup B^c\]
\end{formula}
\guillemotleft{}Ninguno de los dos\guillemotright{} \(=A^c\cap B^c=(A\cup B)^c\). \guillemotleft{}No ambos\guillemotright{} \(=(A\cap B)^c\). \guillemotleft{}Al menos uno\guillemotright{} \(=A\cup B\).

\seccion{3. Frecuencias}
\begin{formula}
\[f_n(A)=\frac{n_A}{n},\qquad 0\le f_n(A)\le1\]
\end{formula}
\(n_A\): frecuencia absoluta (veces que ocurre \(A\) en \(n\) repeticiones). Al aumentar \(n\), la frecuencia relativa tiende a estabilizarse en torno a \(P(A)\) (ley de los grandes números). La probabilidad empírica se estima con datos; la teórica sale del modelo.

\seccion{4. Axiomas de Kolmogórov}
\begin{formula}
1) \(P(A)\ge0\). \quad 2) \(P(\Omega)=1\).\\
3) Si \(A\cap B=\varnothing\): \(P(A\cup B)=P(A)+P(B)\) (y para una sucesión de sucesos incompatibles dos a dos).
\end{formula}

\seccion{5. Propiedades}
\begin{formula}
\[P(A^c)=1-P(A),\qquad P(\varnothing)=0\]
\[P(A\cup B)=P(A)+P(B)-P(A\cap B)\]
\[P(A-B)=P(A)-P(A\cap B)\]
\[A\subset B\ \Rightarrow\ P(A)\le P(B)\]
\end{formula}
\begin{ejemplo}
\(P(A)=0{,}5\), \(P(B)=0{,}4\), \(P(A\cap B)=0{,}2\):
\(P(A\cup B)=0{,}7\); ninguno: \(P(A^c\cap B^c)=1-0{,}7=0{,}3\); solo \(A\): \(0{,}5-0{,}2=0{,}3\).
\end{ejemplo}
\sub{Al menos uno:} mejor por el contrario. Al menos un 6 en tres tiradas: \(1-\bigl(\frac56\bigr)^3=\frac{91}{216}\approx0{,}42\).

\seccion{6. Tablas de contingencia}
Organizan dos características. Las probabilidades se leen dividiendo entre el total (o entre el total de una fila o columna si se condiciona).
\begin{ejemplo}
\[\begin{array}{l|cc|c}
 & H & M & \text{Total}\\\hline
F\ (\text{fuma}) & 20 & 10 & 30\\
F^c & 30 & 40 & 70\\\hline
\text{Total} & 50 & 50 & 100
\end{array}\]
\(P(F)=0{,}3\), \(P(F\cap H)=0{,}2\), \(P(F\mid H)=\frac{20}{50}=0{,}4\), \(P(H\mid F)=\frac{20}{30}=\frac23\).
\end{ejemplo}

\seccion{7. Probabilidad condicionada}
\begin{formula}
\[P(A\mid B)=\frac{P(A\cap B)}{P(B)},\qquad P(B)>0\]
\textbf{Regla del producto:} \(P(A\cap B)=P(B)\,P(A\mid B)\).
\end{formula}
\(P(A\mid B)\) es la probabilidad de \(A\) \textbf{sabiendo que} ha ocurrido \(B\): el espacio se reduce a \(B\).
\begin{ojo}
\(P(A\mid B)\ne P(B\mid A)\) en general: en la tabla, \(0{,}4\ne\frac23\).
\end{ojo}

\seccion{8. Independencia}
\begin{formula}
\(A\) y \(B\) independientes \(\iff P(A\cap B)=P(A)\,P(B)\iff P(A\mid B)=P(A)\).
\end{formula}
En la tabla: \(P(F\cap H)=0{,}2\ne0{,}3\cdot0{,}5=0{,}15\): fumar y sexo \textbf{no} son independientes.
\begin{ojo}
Incompatibles \(\ne\) independientes. Si \(A\) y \(B\) son incompatibles con probabilidades positivas, \(P(A\cap B)=0\ne P(A)P(B)\): son dependientes.
\end{ojo}

\seccion{9. Experimentos compuestos: árboles}
Las probabilidades de cada camino se \textbf{multiplican}; los caminos que forman un suceso se \textbf{suman}.
\begin{ejemplo}
Urna con 3 rojas y 2 azules; se sacan dos bolas.
Con reemplazamiento: \(P(RR)=\frac35\cdot\frac35=\frac9{25}\).
Sin reemplazamiento: \(P(RR)=\frac35\cdot\frac24=\frac3{10}\) (la segunda depende de la primera).
\end{ejemplo}

\seccion{10. Probabilidad total}
Si \(A_1,\dots,A_n\) forman una \textbf{partición} de \(\Omega\) (incompatibles dos a dos y su unión es \(\Omega\)):
\begin{formula}
\[P(B)=\sum_{i=1}^{n}P(A_i)\,P(B\mid A_i)\]
\end{formula}
\begin{center}
\begin{tikzpicture}[grow=right,level distance=1.9cm,level 1/.style={sibling distance=1.5cm},level 2/.style={sibling distance=0.7cm},every node/.style={font=\scriptsize}]
\node {}
  child {node {\(M_2\)}
    child {node {\(D^c\)} edge from parent node[below] {0,95}}
    child {node {\(D\)} edge from parent node[above] {0,05}}
    edge from parent node[below] {0,4}}
  child {node {\(M_1\)}
    child {node {\(D^c\)} edge from parent node[below] {0,98}}
    child {node {\(D\)} edge from parent node[above] {0,02}}
    edge from parent node[above] {0,6}};
\end{tikzpicture}
\end{center}
\begin{ejemplo}
\(M_1\) fabrica el 60\,\% de las piezas (2\,\% defectuosas) y \(M_2\) el 40\,\% (5\,\% defectuosas).
\[P(D)=0{,}6\cdot0{,}02+0{,}4\cdot0{,}05=0{,}012+0{,}02=0{,}032\]
\end{ejemplo}

\seccion{11. Teorema de Bayes}
\begin{formula}
\[P(A_i\mid B)=\frac{P(A_i)\,P(B\mid A_i)}{\sum_{j=1}^{n}P(A_j)\,P(B\mid A_j)}=\frac{P(A_i)\,P(B\mid A_i)}{P(B)}\]
\end{formula}
Invierte la condición: conocida la consecuencia \(B\), \textquestiondown{}qué probabilidad tiene cada causa \(A_i\)?
\begin{ejemplo}
Una pieza es defectuosa. \textquestiondown{}Probabilidad de que venga de \(M_2\)?
\[P(M_2\mid D)=\frac{0{,}4\cdot0{,}05}{0{,}032}=\frac{0{,}02}{0{,}032}=0{,}625\]
\(M_2\) fabrica el 40\,\%, pero aporta el 62,5\,\% de los defectos.
\end{ejemplo}

\seccion{12. Bayes con una prueba diagnóstica}
Ejemplo ficticio: el 1\,\% de una población tiene una enfermedad \(E\). La prueba da positivo en el 95\,\% de los enfermos y en el 5\,\% de los sanos.
\[P(+)=0{,}01\cdot0{,}95+0{,}99\cdot0{,}05=0{,}0095+0{,}0495=0{,}059\]
\[P(E\mid +)=\frac{0{,}0095}{0{,}059}\approx0{,}161\]
Solo un 16\,\% de los positivos están enfermos: como la enfermedad es poco frecuente, la mayoría de los positivos vienen de los muchos sanos. \(P(+\mid E)=0{,}95\) y \(P(E\mid +)\approx0{,}16\) son preguntas distintas.

\seccion{13. De porcentajes a una tabla}
En un centro, el 60\,\% son chicas; el 20\,\% de las chicas y el 30\,\% de los chicos hacen deporte federado. Sobre 100 personas:
\[\begin{array}{l|cc|c}
 & \text{Chica} & \text{Chico} & \text{Total}\\\hline
D & 12 & 12 & 24\\
D^c & 48 & 28 & 76\\\hline
\text{Total} & 60 & 40 & 100
\end{array}\]
\(P(D)=0{,}24\) (probabilidad total) y \(P(\text{Chica}\mid D)=\frac{12}{24}=0{,}5\) (Bayes), sin fórmulas: leyendo la tabla.

\seccion{14. Traducir expresiones}
\begin{tabular}{@{}p{3.2cm}p{4.6cm}@{}}
\textbf{Enunciado} & \textbf{Suceso y probabilidad}\\ \hline
Al menos uno & \(A\cup B\)\\
Ambos & \(A\cap B\)\\
Ninguno & \((A\cup B)^c\): \(1-P(A\cup B)\)\\
Solo \(A\) & \(A\cap B^c\): \(P(A)-P(A\cap B)\)\\
Exactamente uno & \(P(A)+P(B)-2P(A\cap B)\)\\
No ambos & \((A\cap B)^c\): \(1-P(A\cap B)\)\\
\(A\) sabiendo \(B\) & \(P(A\mid B)\)
\end{tabular}

\seccion{15. Estrategia}
\begin{enumerate}
  \item Nombra los sucesos con letras y traduce el enunciado (\guillemotleft{}sabiendo que\guillemotright{} \(\to\) condicionada).
  \item Elige la herramienta: tabla, árbol, Laplace o fórmulas.
  \item \textquestiondown{}Pide una causa conocida la consecuencia? \(\to\) Bayes.
  \item Comprueba: probabilidades entre 0 y 1; las de cada nudo del árbol suman 1.
\end{enumerate}

\begin{ojo}
\begin{itemize}
  \item Sumar \(P(A)+P(B)\) sin restar la intersección.
  \item Confundir \(P(A\cap B)\) con \(P(A\mid B)\).
  \item Multiplicar probabilidades de sucesos que no son independientes sin condicionar.
  \item Olvidar que sin reemplazamiento cambian las probabilidades.
\end{itemize}
\end{ojo}

\end{multicols}
\end{document}
